% Options for packages loaded elsewhere
% Options for packages loaded elsewhere
\PassOptionsToPackage{unicode}{hyperref}
\PassOptionsToPackage{hyphens}{url}
\PassOptionsToPackage{dvipsnames,svgnames,x11names}{xcolor}
%
\documentclass[
  letterpaper,
  DIV=11,
  numbers=noendperiod]{scrartcl}
\usepackage{xcolor}
\usepackage{amsmath,amssymb}
\setcounter{secnumdepth}{5}
\usepackage{iftex}
\ifPDFTeX
  \usepackage[T1]{fontenc}
  \usepackage[utf8]{inputenc}
  \usepackage{textcomp} % provide euro and other symbols
\else % if luatex or xetex
  \usepackage{unicode-math} % this also loads fontspec
  \defaultfontfeatures{Scale=MatchLowercase}
  \defaultfontfeatures[\rmfamily]{Ligatures=TeX,Scale=1}
\fi
\usepackage{lmodern}
\ifPDFTeX\else
  % xetex/luatex font selection
\fi
% Use upquote if available, for straight quotes in verbatim environments
\IfFileExists{upquote.sty}{\usepackage{upquote}}{}
\IfFileExists{microtype.sty}{% use microtype if available
  \usepackage[]{microtype}
  \UseMicrotypeSet[protrusion]{basicmath} % disable protrusion for tt fonts
}{}
\makeatletter
\@ifundefined{KOMAClassName}{% if non-KOMA class
  \IfFileExists{parskip.sty}{%
    \usepackage{parskip}
  }{% else
    \setlength{\parindent}{0pt}
    \setlength{\parskip}{6pt plus 2pt minus 1pt}}
}{% if KOMA class
  \KOMAoptions{parskip=half}}
\makeatother
% Make \paragraph and \subparagraph free-standing
\makeatletter
\ifx\paragraph\undefined\else
  \let\oldparagraph\paragraph
  \renewcommand{\paragraph}{
    \@ifstar
      \xxxParagraphStar
      \xxxParagraphNoStar
  }
  \newcommand{\xxxParagraphStar}[1]{\oldparagraph*{#1}\mbox{}}
  \newcommand{\xxxParagraphNoStar}[1]{\oldparagraph{#1}\mbox{}}
\fi
\ifx\subparagraph\undefined\else
  \let\oldsubparagraph\subparagraph
  \renewcommand{\subparagraph}{
    \@ifstar
      \xxxSubParagraphStar
      \xxxSubParagraphNoStar
  }
  \newcommand{\xxxSubParagraphStar}[1]{\oldsubparagraph*{#1}\mbox{}}
  \newcommand{\xxxSubParagraphNoStar}[1]{\oldsubparagraph{#1}\mbox{}}
\fi
\makeatother


\usepackage{longtable,booktabs,array}
\usepackage{calc} % for calculating minipage widths
% Correct order of tables after \paragraph or \subparagraph
\usepackage{etoolbox}
\makeatletter
\patchcmd\longtable{\par}{\if@noskipsec\mbox{}\fi\par}{}{}
\makeatother
% Allow footnotes in longtable head/foot
\IfFileExists{footnotehyper.sty}{\usepackage{footnotehyper}}{\usepackage{footnote}}
\makesavenoteenv{longtable}
\usepackage{graphicx}
\makeatletter
\newsavebox\pandoc@box
\newcommand*\pandocbounded[1]{% scales image to fit in text height/width
  \sbox\pandoc@box{#1}%
  \Gscale@div\@tempa{\textheight}{\dimexpr\ht\pandoc@box+\dp\pandoc@box\relax}%
  \Gscale@div\@tempb{\linewidth}{\wd\pandoc@box}%
  \ifdim\@tempb\p@<\@tempa\p@\let\@tempa\@tempb\fi% select the smaller of both
  \ifdim\@tempa\p@<\p@\scalebox{\@tempa}{\usebox\pandoc@box}%
  \else\usebox{\pandoc@box}%
  \fi%
}
% Set default figure placement to htbp
\def\fps@figure{htbp}
\makeatother


% definitions for citeproc citations
\NewDocumentCommand\citeproctext{}{}
\NewDocumentCommand\citeproc{mm}{%
  \begingroup\def\citeproctext{#2}\cite{#1}\endgroup}
\makeatletter
 % allow citations to break across lines
 \let\@cite@ofmt\@firstofone
 % avoid brackets around text for \cite:
 \def\@biblabel#1{}
 \def\@cite#1#2{{#1\if@tempswa , #2\fi}}
\makeatother
\newlength{\cslhangindent}
\setlength{\cslhangindent}{1.5em}
\newlength{\csllabelwidth}
\setlength{\csllabelwidth}{3em}
\newenvironment{CSLReferences}[2] % #1 hanging-indent, #2 entry-spacing
 {\begin{list}{}{%
  \setlength{\itemindent}{0pt}
  \setlength{\leftmargin}{0pt}
  \setlength{\parsep}{0pt}
  % turn on hanging indent if param 1 is 1
  \ifodd #1
   \setlength{\leftmargin}{\cslhangindent}
   \setlength{\itemindent}{-1\cslhangindent}
  \fi
  % set entry spacing
  \setlength{\itemsep}{#2\baselineskip}}}
 {\end{list}}
\usepackage{calc}
\newcommand{\CSLBlock}[1]{\hfill\break\parbox[t]{\linewidth}{\strut\ignorespaces#1\strut}}
\newcommand{\CSLLeftMargin}[1]{\parbox[t]{\csllabelwidth}{\strut#1\strut}}
\newcommand{\CSLRightInline}[1]{\parbox[t]{\linewidth - \csllabelwidth}{\strut#1\strut}}
\newcommand{\CSLIndent}[1]{\hspace{\cslhangindent}#1}



\setlength{\emergencystretch}{3em} % prevent overfull lines

\providecommand{\tightlist}{%
  \setlength{\itemsep}{0pt}\setlength{\parskip}{0pt}}



 


\KOMAoption{captions}{tableheading}
\makeatletter
\@ifpackageloaded{caption}{}{\usepackage{caption}}
\AtBeginDocument{%
\ifdefined\contentsname
  \renewcommand*\contentsname{Table of contents}
\else
  \newcommand\contentsname{Table of contents}
\fi
\ifdefined\listfigurename
  \renewcommand*\listfigurename{List of Figures}
\else
  \newcommand\listfigurename{List of Figures}
\fi
\ifdefined\listtablename
  \renewcommand*\listtablename{List of Tables}
\else
  \newcommand\listtablename{List of Tables}
\fi
\ifdefined\figurename
  \renewcommand*\figurename{Figure}
\else
  \newcommand\figurename{Figure}
\fi
\ifdefined\tablename
  \renewcommand*\tablename{Table}
\else
  \newcommand\tablename{Table}
\fi
}
\@ifpackageloaded{float}{}{\usepackage{float}}
\floatstyle{ruled}
\@ifundefined{c@chapter}{\newfloat{codelisting}{h}{lop}}{\newfloat{codelisting}{h}{lop}[chapter]}
\floatname{codelisting}{Listing}
\newcommand*\listoflistings{\listof{codelisting}{List of Listings}}
\makeatother
\makeatletter
\usepackage{pdflscape}
\makeatother
\makeatletter
\makeatother
\makeatletter
\@ifpackageloaded{caption}{}{\usepackage{caption}}
\@ifpackageloaded{subcaption}{}{\usepackage{subcaption}}
\makeatother
\usepackage{bookmark}
\IfFileExists{xurl.sty}{\usepackage{xurl}}{} % add URL line breaks if available
\urlstyle{same}
\hypersetup{
  pdftitle={Evaluating Different Habitat Suitability Modeling for the Eastern Oyster in the Pensacola Bay System},
  pdfauthor={M. Fazio},
  colorlinks=true,
  linkcolor={blue},
  filecolor={Maroon},
  citecolor={Blue},
  urlcolor={Blue},
  pdfcreator={LaTeX via pandoc}}


\title{Evaluating Different Habitat Suitability Modeling for the Eastern
Oyster in the Pensacola Bay System}
\author{M. Fazio}
\date{2026-10-07}
\begin{document}
\maketitle

\renewcommand*\contentsname{Table of contents}
{
\hypersetup{linkcolor=}
\setcounter{tocdepth}{3}
\tableofcontents
}

\begin{center}\rule{0.5\linewidth}{0.5pt}\end{center}

bibliography: \ldots/references.bib

\begin{center}\rule{0.5\linewidth}{0.5pt}\end{center}

\section{Introduction}\label{introduction}

The eastern oyster (\emph{Crassostrea virginica}) are viewed as a
foundational species as they are directly responsible for creating and
providing numerous ecosystem goods and services such as food and habitat
production, shoreline stabilization and protection, water quality
improvement, nutrient cycling, and carbon sequestration (Geselbracht et
al. 2024; Grabowski et al. 2012; Barnes et al. 2007). Additionally, they
provide numerous economic benefits such as supporting oyster fisheries,
recreational harvesting, tourism, along with the associated employment
that goes along with those industries (Geselbracht et al. 2024; Peyre,
Marshall, and Sable 2021). In one study, Grabowski et al. (2012)
estimated the average annual value to range from \$10,325 to \$99,421
per hectare (2011 Dollars).

In recent decades, oyster reefs have experienced substantial declines,
with only an estimated 15\% remaining worldwide, and 50 to 80\% being
lost within the Gulf of (America) alone (Pollack et al. 2012; Peyre,
Marshall, and Sable 2021). The decline is largely attributed to a
combination of anthropogenic activities such us over harvesting, habitat
alteration, sedimentation and water quality decline, disease and
changing hydrology (Peyre, Marshall, and Sable 2021; Geselbracht et al.
2024). In response to these sharp declines, many restoration efforts
have been undertaken with varying strategies and success rates;
recently, planning documents released by the Deepwater Horizon Natural
resource Damage Assessment Trustees called for large-scale development
of reef networks to increase resilience (Geselbracht et al. 2024).

\subsection{Oyster decline in the Pensacola Bay
System}\label{oyster-decline-in-the-pensacola-bay-system}

On such system to experience oyster decline is the Pensacola Bay System
(PBS). The PBS is a large estuarine system located in Escambia and Santa
Rosa counties in northwest Florida. The system is comprised of multiple
waterbody's including Pensacola Bay, Escambia Bay, Blackwater Bay, East
Bay, and Santa Rosa Sound; and is influenced heavily by freshwater input
from the Escambia, BLackwater and Yellow Rivers (Christensen et al.
1997). Historically, the PBS supported extensive oyster reefs that
provided critical habitat for a variety of species and substantial
economic value to the surrounding communities (Christensen et al. 1997).
However, over the past several decades, the extent of oyster habitat
within the PBS has declined significantly due to a combination of
factors including over harvesting, habitat degradation, and water
quality issues. With some estimates indicating a decline of 15,000ha to
just 91.5 ha of oyster habitat between the years of 1883 and 2021
(Geselbracht et al. 2024). This decline has prompted efforts to restore
oyster reefs within the PBS (Pensacola and Perdido Bays Estuary Program
n.d.; The Nature Conservancy 2021), with multiple studies aiming to
identify the most suitable areas within the estuary to focus restoration
efforts (Geselbracht et al. 2024; Christensen et al. 1997; Pensacola and
Perdido Bays Estuary Program n.d.).

With limited resources available for restoration, it is critical to
identify areas where chances of oyster survival and long-term reef
establishment are at their highest, and where poor site selection could
ultimately lead to unsatisfactory results. As with any ecosystem, the
environmental conditions within the PBS vary greatly in time and space,
and such variations can have great impacts on oyster survival and
establishment. Therefore, it is important to understand the
environmental factors that drive oyster occurrence so that scientist and
policy makers can make informed decisions regarding said restoration
site selection (Pollack et al. 2012; Barnes et al. 2007; Geselbracht et
al. 2024).

\subsection{Environmental controls on Eastern oyster
habitat}\label{environmental-controls-on-eastern-oyster-habitat}

Oyster recruitment, survival, and reef establishment are influenced by
complex interactions between numerous physical, chemical, and biological
environmental factors. As such, determining the most influential factors
is a multivariate problem that requires careful consideration of each
factor and its interactions with others (Peyre, Marshall, and Sable
2021; Pollack et al. 2012; Barnes et al. 2007; Christensen et al. 1997).
To help narrow the focus, the most commonly cited environmental factors
were examined in the literature and are summarized here.

Perhaps the most influential and widely reported environmental driver of
oyster suitability is salinity. Representing an important environmental
gradient within estuarine systems that influences physiology,
recruitment, disease, predation, and survival. Other important
characteristics often reported can be seen in
\textbf{?@tbl-environmental-drivers} and include: dissolved oxygen,
temperature, depth, turbidity, and bottom type (substrate) (Peyre,
Marshall, and Sable 2021; Pollack et al. 2012; Barnes et al. 2007;
Geselbracht et al. 2024; Christensen et al. 1997).

\begin{landscape}

\begin{longtable}[]{@{}
  >{\raggedright\arraybackslash}p{(\linewidth - 4\tabcolsep) * \real{0.1236}}
  >{\raggedright\arraybackslash}p{(\linewidth - 4\tabcolsep) * \real{0.5491}}
  >{\raggedright\arraybackslash}p{(\linewidth - 4\tabcolsep) * \real{0.3273}}@{}}
\caption{Common environmental drivers of Eastern Oyster habitat
suitability.}\tabularnewline
\toprule\noalign{}
\begin{minipage}[b]{\linewidth}\raggedright
Factor
\end{minipage} & \begin{minipage}[b]{\linewidth}\raggedright
Description
\end{minipage} & \begin{minipage}[b]{\linewidth}\raggedright
Sources
\end{minipage} \\
\midrule\noalign{}
\endfirsthead
\toprule\noalign{}
\begin{minipage}[b]{\linewidth}\raggedright
Factor
\end{minipage} & \begin{minipage}[b]{\linewidth}\raggedright
Description
\end{minipage} & \begin{minipage}[b]{\linewidth}\raggedright
Sources
\end{minipage} \\
\midrule\noalign{}
\endhead
\bottomrule\noalign{}
\endlastfoot
Salinity & Controls feeding, growth, reproduction, recruitment, disease,
predation, and survival. & (Barnes et al. 2007; Birch et al. 2021;
Christensen et al. 1997; Geselbracht et al. 2024; Peyre, Marshall, and
Sable 2021; Pollack et al. 2012) \\
Temperature & Regulates metabolism, filtration, growth, spawning, and
larval development. & (Barnes et al. 2007; Christensen et al. 1997;
Geselbracht et al. 2024; Peyre, Marshall, and Sable 2021; Pollack et al.
2012) \\
Substrate type & Stable hard substrate supports larval attachment and
reef persistence, whereas soft or rapidly accumulating sediments can
bury or destabilize oysters. & (Barnes et al. 2007; Birch et al. 2021;
Christensen et al. 1997; Geselbracht et al. 2024; Peyre, Marshall, and
Sable 2021) \\
Dissolved oxygen & Adequate oxygen supports respiration, feeding, and
growth. & (Christensen et al. 1997; Geselbracht et al. 2024; Peyre,
Marshall, and Sable 2021; Pollack et al. 2012) \\
Hydrology / water movement & Water movement regulates salinity and
transports larvae, oxygen, and food. & (Barnes et al. 2007; Geselbracht
et al. 2024; Peyre, Marshall, and Sable 2021) \\
Depth / bathymetry & Depth influences exposure, temperature, oxygen
conditions, and restoration feasibility. & (Christensen et al. 1997;
Geselbracht et al. 2024; Pollack et al. 2012) \\
Turbidity / sedimentation & High suspended-sediment loads can interfere
with filtration and settlement. Sediment deposition can smoother oyster
reefs. & (Peyre, Marshall, and Sable 2021; Pollack et al. 2012) \\
Food availability / chlorophyll a & Chlorophyll a is often used as a
proxy for food availability. & (Geselbracht et al. 2024; Peyre,
Marshall, and Sable 2021) \\
\end{longtable}

\end{landscape}

It should be recognized that not all environmental factors are
adequately spatially distributed or represented within the available
datasets, and therefore careful consideration should be given when
selecting the final dataset used for modeling habitat suitability
(Peyre, Marshall, and Sable 2021; Geselbracht et al. 2024; Christensen
et al. 1997).

\subsection{Species Distribution
Modelling}\label{species-distribution-modelling}

Quantifying the relationships between the distribution of species and
the species-environment interactions have a long history in ecological
research, having been described as far back as the late 17th early 18th
century in which researchers observed biological patterns in terms of
their relationships with geographical and/or environmental gradients
(Rushton, Ormerod, and Kerby 2004; Elith and Leathwick 2009). Since
then, the propagation of large-scale spatial datasets, the development
and integration of GIS, increased access to statistical packages and
computational power have allowed for much more development in modeling
complex species-environment interactions (Johnson and Gillingham 2005).
Out of such developments has arisen a technique known as Species
Distribution Modeling (SDM), which can be loosely defined as numerical
tools that combine observations of species occurrence or abundance with
environmental estimates (Elith and Leathwick 2009). According to Elith
and Leathwick (2009), SDMs can be largely classified within several
families of models: expert-driven or heuristic, environmental-envelope,
environmental-similarity, regression-based, Bayesian and hierarchical,
machine-learning and data mining, and ensemble or multimodal.
Additionally, the selection of an SDM approach largely depends on a
multitude of factors such as: the study's objective, quality and
quantity of available data, ecological characteristics of the species,
spatial and temporal scale, whether input data includes absence data or
presence-only data (Elith and Leathwick 2009). With so many choices of
modeling technique and the driving factors behind that choice, deciding
on one model can be quite difficult and the choice usually defaults to a
standard technique depending on the subject matter.

\emph{Habitat Suitability Indices} One popular approach used in many
studies involving the Eastern oyster throughout the gulf region is that
of the expert-driven Habitat Suitability Index (Barnes et al. 2007;
Pollack et al. 2012; Geselbracht et al. 2024; Christensen et al. 1997).
HSIs are typically used in situations where quantitative data on species
distributions is difficult to collect or non-existant, relying on
expert-driven variable selection, meta-analysis, and classification
methods that transform environmental characteristics into standardized
estimates of habitat suitability (Barnes et al. 2007; Pollack et al.
2012; Geselbracht et al. 2024; Johnson and Gillingham 2005).

Within the PBS, HSIs have been used to develop GIS-based and
estuary-specific HSMs that integrate biological, physical, chemical, and
management-related variables in order to identify suitable restoration
areas with the system. In the mid 1990's Christensen et al. (1997)
developed a HSI for the eastern oyster in the PBS to demonstrate the
model frameworks usefulness in illustrating species suitability in a
system wide approach. Geselbracht et al. (2024) then followed and
further expanded on the approach incorporating management factors into
the model, adding a layer of regulatory influence to the model output
versus the purely environmental analysis of the prior. Additionally, the
Pensacola and Perdido Bay Estuary Program recently built on-top of
Geselbracht et al. (2024) by incorporating some of their own data and
refining some of the weighting (Pensacola and Perdido Bays Estuary
Program n.d.).

These previous studies alone highlight the strengths and weakness of the
HSI approaches. Some of the strengths include interoperability,
adaptability, and the integration of ecological knowledge, but the
results entirely rely on expert decisions regarding variable selection,
suitability functions, weighting and scoring, and how to aggregate the
results. All of these present opportunity for differing outputs even if
the same inputs were used by different groups (Geselbracht et al. 2024;
Johnson and Gillingham 2005).

\emph{Maximum Entropy Modelling (MaxEnt)} MaxEnt is a
presence-background SDM framework based on general-purpose machine
learning methods that are useful for making predictions using incomplete
information by relating environmental conditions at known occurrence
locations to the environmental conditions throughout the study area
(Phillips, Anderson, and Schapire 2006; Elith et al. 2011). Some of the
main benefits of MaxEnt is that it uses presence-only data, reducing the
need for the often un-available absence data, it has the ability to
model complex interactions between environmental conditions, and it has
demonstrated high predictive performance while limiting over fitting
through regularization (Phillips et al. 2009; Merow, Smith, and Silander
2013; Elith et al. 2011). An added benefit to this particular study is
the fact that MaxEnt is tailored towards presence-only species
occurrence data, making it highly suitable for the data set used here in
which the FWC-FWRI (2026) datasets explicitly collected occurrence data
without including absence locations.

However, the use of such algorithms should not be made blindly as many
issues can arise if careful precautions are not taken. Due to the nature
of ecological many biases in the data samples can arise such as spatial
and temporal autocorrelation, group dependence structures, hierarchical
structures and multicollinearity leading to overfitting model parameters
and over-optimistic prediction confidence (Dormann et al. 2013; Kass et
al. 2021; Roberts et al. 2017). Thankfully, significant work has been
put into this through cross-validation techniques and the development of
model evaluation software packages such as ENMeval 2.0 implemented
through the R statistical language (Kass et al. 2021; Roberts et al.
2017)

\subsection{Research rationale, objectives, and
hypotheses}\label{research-rationale-objectives-and-hypotheses}

In an effort to demonstrate how a MaxEnt SDM can be used as a
complimentary tool to help researchers and practitioners leading
ecological restoration activities, this study aims to develop an
empirically fitted MaxEnt habitat-suitability model using observed
oyster occurrence data and known environmental conditions. The primary
objective will be to develop a spatially explicit SDM representing
relative Eastern oyster habitat suitabiliy throughout the PBS. Next, the
resulting suitability surface will be evaluated agains the existing
Pensacola Bay HSI developed by Geselbracht et al. (2024) and iterated
upon by Pensacola and Perdido Bays Estuary Program (n.d.).

To achieve this, we will identify the environmental predictors that are
commonly associated with oyster habitat success,

\begin{enumerate}
\def\labelenumi{\arabic{enumi}.}
\setcounter{enumi}{3}
\tightlist
\item
  Supporting objectives.

  \begin{itemize}
  \tightlist
  \item
    Identify environmental predictors associated with observed oyster
    distribution.
  \item
    Evaluate MaxEnt predictive performance using spatially separated
    validation data.
  \item
    Identify geographic areas where MaxEnt and the HSI agree or
    disagree.
  \item
    Evaluate whether MaxEnt provides a reproducible empirical complement
    to HSI-based restoration planning.
  \item
    Citations: (Johnson and Gillingham 2005; Geselbracht et al. 2024;
    Merow, Smith, and Silander 2013)
  \end{itemize}
\item
  Research hypothesis.

  \begin{itemize}
  \tightlist
  \item
    The MaxEnt model is expected to identify patterns of habitat
    suitability consistent with established Eastern oyster habitat
    requirements.
  \item
    Measurable agreement is expected between MaxEnt and the existing
    HSI, with spatial disagreement arising from differences between
    empirically fitted environmental relationships and predefined HSI
    suitability relationships.
  \item
    Citations: (Johnson and Gillingham 2005; Geselbracht et al. 2024)
  \end{itemize}
\end{enumerate}

\section{Methods}\label{methods}

\emph{Study design and team status}

This capstone is being completed as an individual project. As such, all
data acquisition, processing, modeling, and interpretations will be
produced soley by the author.

All data analysis and processing will be performed using the R
programming language (R Core Team 2024). Environmental observation data
will be acquired from the EPA's Water Quality Portal via the
dataRetrieval r-package (USGS 2026; DeCicco et al. 2026). Data QA/QC and
preprocessing will be completed utilizing the Tidyverse package and any
spatially aware data processing will be conducted using the package sf
(Wickham et al. 2019; Pebesma 2018). Environmental data will be
processed to create continuous surface rasters for model input using the

\subsubsection{Study Area}\label{study-area}

The study area for this project is the Pensacola Bay System which
includes portions of Pensacola Bay, Escambia Bay, East Bay and
Blackwater Bay. Boundaries of this study area were achieved by accessing
Florida Department of Protections (FDEP) Waterbod IDs feature layer
hosted on ArcGIS Online (FDEP 2005). ArcGIS allows direct access to
hosted feature services via several formats. To avoid any potential
compatibility issues, the GeoJSON format was chosen as this format is
non-proprietary. Once retrieved, the dataset was filtered to include
only relevant WBIDs (WBID = 548) and then combined to create a singular
study area polygon. Creating the study area polygon is an important
first step as it will be needed to extract relevant environmental data
from much larger datasets in the coming data gathering.

\begin{figure}[H]

{\centering \pandocbounded{\includegraphics[keepaspectratio]{"../4-tables-and-figures/fig-1-study_area.png"}}

}

\caption{Study Area}

\end{figure}%

\begin{enumerate}
\def\labelenumi{\arabic{enumi}.}
\setcounter{enumi}{1}
\item
  Describe the overall analytical workflow.

  \begin{itemize}
  \tightlist
  \item
    Environmental and occurrence data acquisition.
  \item
    Data QA/QC and preprocessing.
  \item
    Creation of continuous environmental predictor surfaces.
  \item
    Oyster occurrence preparation.
  \item
    Background sampling.
  \item
    Predictor screening.
  \item
    MaxEnt tuning and fitting (Phillips, Anderson, and Schapire 2006;
    Merow, Smith, and Silander 2013; Kass et al. 2021).
  \item
    Spatial model validation (Roberts et al. 2017; Valavi et al. 2019).
  \item
    HSI comparison (Johnson and Gillingham 2005; Geselbracht et al.
    2024).
  \end{itemize}
\item
  \textbf{Study area}

  \begin{enumerate}
  \def\labelenumii{\arabic{enumii}.}
  \tightlist
  \item
    Define the Pensacola Bay System modeling extent.

    \begin{itemize}
    \tightlist
    \item
      Delineate the study area using selected Florida Department of
      Environmental Protection Waterbody Identification (WBID) polygons
      (FDEP 2005).
    \item
      Describe the primary embayments included within the analysis and
      their relationship to the Pensacola Bay System (Geselbracht et al.
      2024).
    \end{itemize}
  \item
    Restrict the modeling domain to aquatic habitat.

    \begin{itemize}
    \tightlist
    \item
      Background and prediction locations should represent environments
      potentially available to Eastern oysters (Merow, Smith, and
      Silander 2013).
    \item
      Terrestrial areas should therefore be excluded from the
      accessible/background modeling domain (Phillips et al. 2009;
      Merow, Smith, and Silander 2013).
    \end{itemize}
  \item
    Maintain spatial compatibility with the existing Pensacola Bay HSI
    where possible.

    \begin{itemize}
    \tightlist
    \item
      A comparable geographic extent will support subsequent spatial
      comparison with the HSI developed for PBS (Geselbracht et al.
      2024).
    \end{itemize}
  \item
    Develop a study-area map.

    \begin{itemize}
    \tightlist
    \item
      PBS boundary (FDEP 2005).
    \item
      Major embayments (Geselbracht et al. 2024).
    \item
      Oyster occurrence locations (FWC-FWRI 2026).
    \item
      Environmental monitoring locations (USGS 2026).
    \end{itemize}
  \end{enumerate}
\item
  \textbf{Data acquisition and dataset inventory}

  \begin{enumerate}
  \def\labelenumii{\arabic{enumii}.}
  \tightlist
  \item
    Eastern oyster occurrence data.

    \begin{itemize}
    \tightlist
    \item
      Source oyster occurrence data from the FWC-FWRI Oyster Beds in
      Florida dataset (FWC-FWRI 2026).
    \item
      Record dataset version and acquisition date (FWC-FWRI 2026).
    \item
      Document original geometry type.
    \item
      Document number of source features and number retained within the
      study area.
    \item
      Document temporal and source-data information where available
      (FWC-FWRI 2026).
    \end{itemize}
  \item
    Water-quality observations.

    \begin{itemize}
    \tightlist
    \item
      Retrieve water-quality observations from the Water Quality Portal
      (USGS 2026).
    \item
      Use the \texttt{dataRetrieval} package and associated WQP services
      for programmatic data acquisition (DeCicco et al. 2026).
    \item
      Restrict queries to relevant counties, estuarine monitoring
      locations, selected parameters, and the study period.
    \end{itemize}
  \item
    Document water-quality dataset size and temporal coverage.

    \begin{itemize}
    \tightlist
    \item
      Dissolved oxygen: approximately 145 stations and 3,072
      observations (USGS 2026).
    \item
      Salinity: approximately 135 stations and 3,401 observations (USGS
      2026).
    \item
      Temperature: approximately 145 stations and 3,442 observations
      (USGS 2026).
    \item
      Turbidity: approximately 122 stations and 1,563 observations (USGS
      2026).
    \item
      Chlorophyll-a: approximately 48 stations and 754 observations
      (USGS 2026).
    \item
      Current temporal coverage is approximately 2005--2024 (USGS 2026).
    \end{itemize}
  \item
    Evaluate candidate predictors based on spatial coverage.

    \begin{itemize}
    \tightlist
    \item
      Chlorophyll-a is expected to be excluded from interpolation
      because available monitoring stations do not provide adequate
      system-wide spatial coverage.
    \item
      Salinity, dissolved oxygen, temperature, and turbidity currently
      provide substantially stronger monitoring coverage.
    \item
      These variables have previously been identified as relevant
      environmental characteristics for Eastern oyster habitat and
      restoration modeling (Peyre, Marshall, and Sable 2021; Barnes et
      al. 2007; Pollack et al. 2012; Geselbracht et al. 2024).
    \end{itemize}
  \item
    Bathymetry.

    \begin{itemize}
    \tightlist
    \item
      Obtain Pensacola Bay estuarine bathymetry from NOAA (National
      Centers for Environmental Information 2018).
    \item
      Document raster resolution, vertical datum, geographic extent, and
      preprocessing requirements (National Centers for Environmental
      Information 2018).
    \end{itemize}
  \item
    Substrate or bottom type.

    \begin{itemize}
    \tightlist
    \item
      Identify and acquire a suitable spatial dataset describing bottom
      characteristics.
    \item
      Document spatial resolution, substrate classes, temporal source,
      and geographic coverage.
    \item
      Suitable substrate is an important constraint on Eastern oyster
      settlement and reef development (Barnes et al. 2007; Geselbracht
      et al. 2024).
    \item
      \textbf{Add a citation for the final substrate dataset once
      selected.}
    \end{itemize}
  \item
    Existing Pensacola Bay HSI.

    \begin{itemize}
    \tightlist
    \item
      Obtain the original GIS/raster output associated with the
      Pensacola Bay HSI developed by Geselbracht et al. (Geselbracht et
      al. 2024).
    \item
      Document raster score range, coordinate reference system, extent,
      resolution, and NoData handling.
    \item
      The HSI will provide the comparison surface for evaluating
      agreement with the MaxEnt model (Geselbracht et al. 2024).
    \item
      \textbf{Add a separate dataset citation once the original HSI GIS
      product is acquired.}
    \end{itemize}
  \end{enumerate}
\item
  \textbf{Water-quality preprocessing}

  \begin{enumerate}
  \def\labelenumii{\arabic{enumii}.}
  \tightlist
  \item
    Standardize measurement characteristics and units.

    \begin{itemize}
    \tightlist
    \item
      Avoid combining fundamentally different measurement types that
      share similar WQP characteristic names (USGS 2026).
    \item
      Retain comparable measurement units and analytical forms for each
      environmental predictor (USGS 2026; DeCicco et al. 2026).
    \end{itemize}
  \item
    Restrict observations to warm-season conditions.

    \begin{itemize}
    \tightlist
    \item
      Define the warm season as May through October.
    \item
      Use seasonal filtering to characterize environmental conditions
      during the portion of the year most relevant to warm-season oyster
      habitat conditions.
    \item
      Seasonal environmental summaries have also been incorporated into
      previous Pensacola Bay oyster suitability modeling (Geselbracht et
      al. 2024).
    \end{itemize}
  \item
    Generate station-level environmental summaries.

    \begin{itemize}
    \tightlist
    \item
      Calculate the warm-season mean for each monitoring location.
    \item
      Calculate the standard deviation for each monitoring location.
    \item
      Record total observation count.
    \item
      Record the number of unique years represented.
    \item
      Record first and last sample dates.
    \item
      Consider calculating site-year warm-season means before
      calculating long-term site means to prevent heavily sampled
      individual years from disproportionately influencing long-term
      conditions.
    \end{itemize}
  \item
    Perform QA/QC.

    \begin{itemize}
    \tightlist
    \item
      Identify potential unit inconsistencies.
    \item
      Investigate implausible or extreme values.
    \item
      Document removal or correction criteria.
    \item
      Retain original raw data separately from processed datasets.
    \end{itemize}
  \item
    Convert monitoring data to spatial features.

    \begin{itemize}
    \tightlist
    \item
      Use the \texttt{sf} package for spatial-vector processing (Pebesma
      2018).
    \item
      Project monitoring locations to NAD 1983 (2011) StatePlane Florida
      North FIPS 0903 in meters.
    \end{itemize}
  \end{enumerate}
\item
  \textbf{Generation of continuous environmental predictor surfaces}

  \begin{enumerate}
  \def\labelenumii{\arabic{enumii}.}
  \tightlist
  \item
    Assess monitoring-station spatial coverage prior to interpolation.

    \begin{itemize}
    \tightlist
    \item
      Map station distributions for each environmental parameter.
    \item
      Identify spatial clustering.
    \item
      Identify portions of PBS lacking observations.
    \item
      Exclude variables where development of a continuous surface would
      consist primarily of spatial extrapolation.
    \end{itemize}
  \item
    Account for coastal barriers and complex estuarine geometry.

    \begin{itemize}
    \tightlist
    \item
      Straight-line Euclidean distance can incorrectly represent spatial
      relationships between monitoring locations separated by
      peninsulas, islands, and other land barriers (Stachelek and Madden
      2015).
    \item
      Barrier-aware interpolation will therefore be evaluated for the
      PBS environmental surfaces (Stachelek and Madden 2015).
    \end{itemize}
  \item
    Evaluate inverse path distance weighting (IPDW) as a candidate
    interpolation method.

    \begin{itemize}
    \tightlist
    \item
      IPDW uses aquatic-path distance rather than straight-line
      Euclidean distance when determining spatial relationships among
      observations (Stachelek and Madden 2015).
    \item
      This approach was developed for coastal water-quality
      interpolation where land barriers influence spatial connectivity
      (Stachelek and Madden 2015).
    \end{itemize}
  \item
    Compare candidate interpolation approaches.

    \begin{itemize}
    \tightlist
    \item
      Evaluate inverse path distance weighting (IPDW) (Stachelek and
      Madden 2015).
    \item
      Evaluate conventional inverse distance weighting (IDW).
    \item
      Evaluate ordinary kriging.
    \item
      Use spatially withheld monitoring locations to evaluate predictive
      performance and reduce bias associated with spatial dependence
      (Roberts et al. 2017; Valavi et al. 2019).
    \item
      Compare root mean square error (RMSE), mean absolute error (MAE),
      and prediction bias among interpolation methods.
    \item
      Select the final interpolation method based on predictive
      performance and ecological appropriateness.
    \end{itemize}
  \item
    Produce final environmental predictor rasters.

    \begin{itemize}
    \tightlist
    \item
      Use a common projected coordinate reference system.
    \item
      Use a common cell size.
    \item
      Use a common raster origin and alignment.
    \item
      Mask predictions to aquatic habitat.
    \item
      Document the final raster resolution.
    \item
      Maintain spatial compatibility with the existing PBS
      habitat-suitability framework where appropriate (Geselbracht et
      al. 2024).
    \end{itemize}
  \end{enumerate}
\item
  \textbf{Oyster occurrence preparation}

  \begin{enumerate}
  \def\labelenumii{\arabic{enumii}.}
  \tightlist
  \item
    Clip FWC oyster-bed features to the study area.

    \begin{itemize}
    \tightlist
    \item
      Remove oyster occurrences outside the PBS modeling extent
      (FWC-FWRI 2026).
    \end{itemize}
  \item
    Convert mapped oyster-bed polygons into representative MaxEnt
    presence observations.

    \begin{itemize}
    \tightlist
    \item
      Use a reproducible procedure that avoids disproportionately
      representing large or densely mapped oyster polygons (Phillips et
      al. 2009; Merow, Smith, and Silander 2013).
    \item
      Document the initial number of oyster polygons and final number of
      presence observations.
    \end{itemize}
  \item
    Evaluate spatial clustering of presence observations.

    \begin{itemize}
    \tightlist
    \item
      Apply spatial thinning where necessary to reduce clustering and
      sampling bias (Phillips et al. 2009).
    \item
      Base thinning decisions on spatial dependence and sampling
      structure rather than arbitrary data reduction (Phillips et al.
      2009; Merow, Smith, and Silander 2013).
    \end{itemize}
  \end{enumerate}
\item
  \textbf{Background sampling}

  \begin{enumerate}
  \def\labelenumii{\arabic{enumii}.}
  \tightlist
  \item
    Define the accessible/background modeling region.

    \begin{itemize}
    \tightlist
    \item
      Restrict background sampling to aquatic habitat within the
      Pensacola Bay System (Merow, Smith, and Silander 2013).
    \item
      Require complete environmental predictor coverage within the
      background domain.
    \item
      Exclude terrestrial areas and locations outside the modeled
      environmental domain (Phillips et al. 2009; Merow, Smith, and
      Silander 2013).
    \end{itemize}
  \item
    Generate MaxEnt background locations.

    \begin{itemize}
    \tightlist
    \item
      Document the final number of background points.
    \item
      Sample background locations from the defined accessible region
      rather than from an arbitrary rectangular geographic extent
      (Phillips et al. 2009; Merow, Smith, and Silander 2013).
    \end{itemize}
  \item
    Address occurrence-sampling bias.

    \begin{itemize}
    \tightlist
    \item
      Evaluate whether spatial thinning, bias weighting, or restricted
      background sampling is necessary (Phillips et al. 2009).
    \end{itemize}
  \end{enumerate}
\item
  \textbf{Predictor screening}

  \begin{enumerate}
  \def\labelenumii{\arabic{enumii}.}
  \tightlist
  \item
    Assemble the final candidate predictor stack.

    \begin{itemize}
    \tightlist
    \item
      Salinity (Peyre, Marshall, and Sable 2021; Barnes et al. 2007;
      Pollack et al. 2012; Geselbracht et al. 2024).
    \item
      Dissolved oxygen (Peyre, Marshall, and Sable 2021; Pollack et al.
      2012; Geselbracht et al. 2024).
    \item
      Temperature (Peyre, Marshall, and Sable 2021; Barnes et al. 2007;
      Pollack et al. 2012).
    \item
      Turbidity (Peyre, Marshall, and Sable 2021; Pollack et al. 2012).
    \item
      Bathymetry/depth (Peyre, Marshall, and Sable 2021; Barnes et al.
      2007; Pollack et al. 2012).
    \item
      Substrate/bottom type, if suitable data are obtained (Peyre,
      Marshall, and Sable 2021; Barnes et al. 2007; Geselbracht et al.
      2024).
    \end{itemize}
  \item
    Evaluate correlation and redundancy among predictors.

    \begin{itemize}
    \tightlist
    \item
      Calculate pairwise correlations among candidate environmental
      variables (Dormann et al. 2013).
    \item
      Identify highly correlated or ecologically redundant predictors
      (Dormann et al. 2013).
    \item
      Retain predictors based on ecological relevance, data quality, and
      interpretability while limiting problematic collinearity (Dormann
      et al. 2013).
    \end{itemize}
  \end{enumerate}
\item
  \textbf{MaxEnt model fitting and tuning}

  \begin{enumerate}
  \def\labelenumii{\arabic{enumii}.}
  \item
    Fit MaxEnt using oyster occurrence locations, background samples,
    and the finalized environmental predictor stack (Phillips, Anderson,
    and Schapire 2006; Elith et al. 2011).
  \item
    Tune model complexity rather than relying exclusively on default
    settings.

    \begin{itemize}
    \tightlist
    \item
      Evaluate multiple regularization multipliers (Merow, Smith, and
      Silander 2013; Kass et al. 2021).
    \item
      Evaluate multiple feature-class combinations (Merow, Smith, and
      Silander 2013; Kass et al. 2021).
    \end{itemize}
  \item
    Use the \texttt{ENMeval} framework for model tuning and evaluation
    (Kass et al. 2021).
  \item
    Select the final MaxEnt configuration.

    \begin{itemize}
    \tightlist
    \item
      Consider predictive performance (Kass et al. 2021).
    \item
      Consider omission rates (Kass et al. 2021).
    \item
      Consider model complexity and overfitting (Merow, Smith, and
      Silander 2013; Kass et al. 2021).
    \end{itemize}
  \item
    Generate the final continuous habitat-suitability surface.

    \begin{itemize}
    \tightlist
    \item
      Retain continuous predictions for the primary spatial comparison
      with the HSI (Elith et al. 2011; Merow, Smith, and Silander 2013).
    \end{itemize}
  \item
    Evaluate environmental-variable influence.

    \begin{itemize}
    \tightlist
    \item
      Examine permutation importance (Merow, Smith, and Silander 2013).
    \item
      Examine model response curves (Elith et al. 2011; Merow, Smith,
      and Silander 2013).
    \item
      Evaluate other supported measures of predictor contribution
      (Merow, Smith, and Silander 2013).
    \item
      Interpret these associations as predictive relationships rather
      than direct evidence of ecological causation (Elith et al. 2011).
    \end{itemize}
  \end{enumerate}
\item
  \textbf{Spatial model validation}

  \begin{enumerate}
  \def\labelenumii{\arabic{enumii}.}
  \tightlist
  \item
    Use spatial rather than purely random cross-validation.

    \begin{itemize}
    \tightlist
    \item
      Nearby oyster occurrences may experience highly similar
      environmental conditions, violating assumptions of independence
      (Roberts et al. 2017).
    \item
      Random training/test partitions may therefore produce overly
      optimistic estimates of predictive performance (Roberts et al.
      2017; Valavi et al. 2019).
    \end{itemize}
  \item
    Divide oyster occurrences into spatially separated validation folds.

    \begin{itemize}
    \tightlist
    \item
      Select spatial-block dimensions based on observed spatial
      structure rather than arbitrary distances (Valavi et al. 2019).
    \end{itemize}
  \item
    Evaluate candidate models against withheld spatial folds.

    \begin{itemize}
    \tightlist
    \item
      Report appropriate presence-background model-performance metrics
      (Phillips, Anderson, and Schapire 2006; Kass et al. 2021).
    \item
      Evaluate omission rates and AUC where appropriate (Phillips,
      Anderson, and Schapire 2006; Kass et al. 2021).
    \item
      Report variation in model performance among spatial folds (Valavi
      et al. 2019; Kass et al. 2021).
    \end{itemize}
  \end{enumerate}
\item
  \textbf{Comparison with the existing Habitat Suitability Index}

  \begin{enumerate}
  \def\labelenumii{\arabic{enumii}.}
  \item
    Harmonize the MaxEnt and HSI raster surfaces.

    \begin{itemize}
    \tightlist
    \item
      Match geographic extent to permit spatial comparison (Johnson and
      Gillingham 2005; Geselbracht et al. 2024).
    \item
      Match coordinate reference systems.
    \item
      Match cell size and raster alignment.
    \item
      Ensure suitability scores are represented on directly comparable
      scales where possible (Johnson and Gillingham 2005).
    \end{itemize}
  \item
    Compare continuous habitat-suitability predictions.

    \begin{itemize}
    \tightlist
    \item
      Calculate cell-by-cell correlation between HSI and MaxEnt
      predictions (Johnson and Gillingham 2005).
    \item
      Consider Pearson and/or Spearman correlation depending on the
      characteristics of the model-output distributions (Johnson and
      Gillingham 2005).
    \end{itemize}
  \item
    Generate a spatial difference surface (Johnson and Gillingham 2005).

    \[
    D_i = S_{MaxEnt,i} - S_{HSI,i}
    \]

    \begin{itemize}
    \tightlist
    \item
      Values near zero represent relatively strong agreement between the
      two suitability estimates.
    \item
      Positive values identify locations where MaxEnt predicts
      relatively greater suitability.
    \item
      Negative values identify locations where the HSI predicts
      relatively greater suitability.
    \end{itemize}
  \item
    Evaluate categorical agreement if suitability classes are created.

    \begin{itemize}
    \tightlist
    \item
      Identify locations where both models predict high suitability
      (Johnson and Gillingham 2005).
    \item
      Identify locations where both models predict low suitability
      (Johnson and Gillingham 2005).
    \item
      Identify locations where MaxEnt predicts high suitability while
      the HSI predicts low suitability (Johnson and Gillingham 2005).
    \item
      Identify locations where the HSI predicts high suitability while
      MaxEnt predicts low suitability (Johnson and Gillingham 2005).
    \item
      Consider an agreement statistic such as Kappa if categorical
      suitability classes are used (Johnson and Gillingham 2005).
    \end{itemize}
  \item
    Interpret disagreement between models.

    \begin{itemize}
    \tightlist
    \item
      Determine whether disagreement is geographically concentrated
      (Johnson and Gillingham 2005).
    \item
      Evaluate whether disagreement corresponds with particular
      environmental predictors.
    \item
      Consider differences in data availability and modeling assumptions
      between MaxEnt and the HSI (Johnson and Gillingham 2005;
      Geselbracht et al. 2024).
    \item
      Avoid treating either modeling approach as automatically
      representing ground truth (Johnson and Gillingham 2005).
    \end{itemize}
  \end{enumerate}
\item
  \textbf{Software and reproducibility}

  \begin{enumerate}
  \def\labelenumii{\arabic{enumii}.}
  \item
    Conduct data processing, statistical analysis, and modeling in R (R
    Core Team 2024).
  \item
    Use scripted workflows for data transformation and spatial
    processing.

    \begin{itemize}
    \tightlist
    \item
      Use \texttt{tidyverse} for data manipulation (Wickham et al.
      2019).
    \item
      Use \texttt{sf} for vector spatial processing (Pebesma 2018).
    \end{itemize}
  \item
    Maintain reproducibility throughout the workflow.

    \begin{itemize}
    \tightlist
    \item
      Record package versions.
    \item
      Set random seeds where applicable.
    \item
      Record data-acquisition dates for dynamically updated datasets
      (DeCicco et al. 2026; USGS 2026; FWC-FWRI 2026).
    \item
      Document model settings and tuning parameters (Merow, Smith, and
      Silander 2013; Kass et al. 2021).
    \item
      Preserve raw data separately from processed analytical datasets.
    \end{itemize}
  \end{enumerate}
\end{enumerate}

\section{Results}\label{results}

In progress\ldots{}

\section{Discussion}\label{discussion}

In progress\ldots{}

\section{Conclusion}\label{conclusion}

In progress\ldots{}

\section{References}\label{references}

\phantomsection\label{refs}
\begin{CSLReferences}{1}{0}
\bibitem[\citeproctext]{ref-Barnes2007}
Barnes, T. K., A. K. Volety, K. Chartier, F. J. Mazzotti, and L.
Pearlstine. 2007. {``A Habitat Suitability Index Model for the Eastern
Oyster ({Crassostrea virginica}), a Tool for Restoration of the
Caloosahatchee Estuary, Florida.''} \emph{Journal of Shellfish Research}
26 (4): 949--59.
\url{https://doi.org/10.2983/0730-8000(2007)26\%5B949:AHSIMF\%5D2.0.CO;2}.

\bibitem[\citeproctext]{ref-Birch2021}
Birch, A., R. Brumbaugh, B. DeAngelis, L. Geselbracht, A. Graves, J.
Blair, and R. Jones. 2021. {``Oyster Fisheries and Habitat Management
Plan for the Pensacola Bay System.''} The Nature Conservancy.
\url{https://www.nature.org/content/dam/tnc/nature/en/documents/PBS_FloridaOyster-Fisheries-Habitat-Mgt-Plan_18May2021_Final.pdf}.

\bibitem[\citeproctext]{ref-Christensen1997}
Christensen, J D, T A Battista, M E Monaco, and C J Klein. 1997.
{``Habitat Suitability Index Modeling and GIS Technology to Support
Habitat Management: Pensacola Bay, Florida Case Study.''} U.S. National
Oceanic; Atmospheric Administration.
\url{https://repository.library.noaa.gov/view/noaa/55683}.

\bibitem[\citeproctext]{ref-dataRetrieval2026}
DeCicco, Laura, Robert Hirsch, David Lorenz, Jordan Read, Jordan Walker,
Lindsay Platt, David Watkins, et al. 2026. {``dataRetrieval: R Packages
for Discovering and Retrieving Water Data Available from u.s. Federal
Hydrologic Web Services.''} Reston, VA: U.S. Geological Survey; U.S.
Geological Survey. \url{https://doi.org/10.5066/P9X4L3GE}.

\bibitem[\citeproctext]{ref-Dormann2013}
Dormann, Carsten F., Jane Elith, Sven Bacher, Carsten Buchmann, Gudrun
Carl, Gabriel Carré, Jaime R.García Marquéz, et al. 2013.
{``Collinearity: A Review of Methods to Deal with It and a Simulation
Study Evaluating Their Performance.''} \emph{Ecography} 36: 27--46.
\url{https://doi.org/10.1111/j.1600-0587.2012.07348.x}.

\bibitem[\citeproctext]{ref-Elith2009}
Elith, Jane, and John R. Leathwick. 2009. {``Species Distribution
Models: Ecological Explanation and Prediction Across Space and Time.''}
\emph{Annual Review of Ecology, Evolution, and Systematics} 40
(December): 677--97.
\url{https://doi.org/10.1146/annurev.ecolsys.110308.120159}.

\bibitem[\citeproctext]{ref-Elith2011}
Elith, Jane, Steven J. Phillips, Trevor Hastie, Miroslav Dudík, Yung En
Chee, and Colin J. Yates. 2011. {``A Statistical Explanation of MaxEnt
for Ecologists.''} \emph{Diversity and Distributions} 17 (January):
43--57. \url{https://doi.org/10.1111/j.1472-4642.2010.00725.x}.

\bibitem[\citeproctext]{ref-FDEP2005WBIDs}
FDEP. 2005. {``Waterbody IDs (WBIDs).''} Geospatial Open Data; Florida
Department of Environmental Protection.
\url{https://geodata.dep.state.fl.us/datasets/FDEP::waterbody-ids-wbids/about}.

\bibitem[\citeproctext]{ref-FWC2026Oysters}
FWC-FWRI. 2026. {``Oyster Beds in Florida.''} Geospatial Open Data;
Florida Fish; Wildlife Conservation Commission, Fish; Wildlife Research
Institute.
\url{https://geodata.myfwc.com/datasets/myfwc::oyster-beds-in-florida/about}.

\bibitem[\citeproctext]{ref-Geselbracht2024}
Geselbracht, L., M. Johnston, B. M. DeAngelis, and A. Birch. 2024.
{``Estuary-Specific and Adaptive Habitat Suitability Index Model for the
Eastern Oyster Crassostrea Virginica in the Pensacola Bay System,
Florida, USA.''} \emph{Coastal Management} 52: 17--34.
\url{https://doi.org/10.1080/08920753.2024.2335129}.

\bibitem[\citeproctext]{ref-Grabowski2012}
Grabowski, Jonathan H., Robert D. Brumbaugh, Robert F. Conrad, Andrew G.
Keeler, James J. Opaluch, Charles H. Peterson, Michael F. Piehler, Sean
P. Powers, and Ashley R. Smyth. 2012. {``Economic Valuation of Ecosystem
Services Provided by Oyster Reefs.''} \emph{BioScience} 62 (October):
900--909. \url{https://doi.org/10.1525/bio.2012.62.10.10}.

\bibitem[\citeproctext]{ref-Johnson2005}
Johnson, Chris J., and Michael P. Gillingham. 2005. {``An Evaluation of
Mapped Species Distribution Models Used for Conservation Planning.''}
\emph{Environmental Conservation} 32 (June): 117--28.
\url{https://doi.org/10.1017/S0376892905002171}.

\bibitem[\citeproctext]{ref-Kass2021}
Kass, Jamie M., Robert Muscarella, Peter J. Galante, Corentin L. Bohl,
Gonzalo E. Pinilla-Buitrago, Robert A. Boria, Mariano Soley-Guardia, and
Robert P. Anderson. 2021. {``ENMeval 2.0: Redesigned for Customizable
and Reproducible Modeling of Species' Niches and Distributions.''}
\emph{Methods in Ecology and Evolution} 12 (September): 1602--8.
\url{https://doi.org/10.1111/2041-210X.13628}.

\bibitem[\citeproctext]{ref-Merow2013}
Merow, Cory, Matthew J. Smith, and John A. Silander. 2013. {``A
Practical Guide to MaxEnt for Modeling Species' Distributions: What It
Does, and Why Inputs and Settings Matter.''} \emph{Ecography} 36
(October): 1058--69.
\url{https://doi.org/10.1111/j.1600-0587.2013.07872.x}.

\bibitem[\citeproctext]{ref-NOAA_2018}
National Centers for Environmental Information. 2018. {``NOAA NOS
Estuarine Bathymetry - Pensacola Bay (G130).''}
\url{https://doi.org/10.7289/V5XK8CWK}.

\bibitem[\citeproctext]{ref-sf}
Pebesma, Edzer. 2018. {``{Simple Features for R: Standardized Support
for Spatial Vector Data}.''} \emph{{The R Journal}} 10 (1): 439--46.
\url{https://doi.org/10.32614/RJ-2018-009}.

\bibitem[\citeproctext]{ref-PPBEP_OysterRestoration}
Pensacola and Perdido Bays Estuary Program. n.d. {``Pensacola Bay Oyster
Restoration Initiative.''}
\url{https://www.ppbep.org/what-we-do/restoration/pensacola-bay-oyster-restoration-initiative}.

\bibitem[\citeproctext]{ref-LaPeyre2021}
Peyre, Megan K. La, Danielle A. Marshall, and Shaye E. Sable. 2021.
{``Oyster Model Inventory: Identifying Critical Data and Modeling
Approaches to Support Restoration of Oyster Reefs in Coastal u.s. Gulf
of Mexico Waters.''} \emph{Diversity and Distributions}. Vol. 17. U.S.
Geological Survey. \url{https://doi.org/10.3133/ofr20211063}.

\bibitem[\citeproctext]{ref-Phillips2006}
Phillips, Steven J., Robert P. Anderson, and Robert E. Schapire. 2006.
{``Maximum Entropy Modeling of Species Geographic Distributions.''}
\emph{Ecological Modelling} 190 (January): 231--59.
\url{https://doi.org/10.1016/j.ecolmodel.2005.03.026}.

\bibitem[\citeproctext]{ref-Phillips2009}
Phillips, Steven J., Miroslav Dudík, Jane Elith, Catherine H. Graham,
Anthony Lehmann, John Leathwick, and Simon Ferrier. 2009. {``Sample
Selection Bias and Presence-Only Distribution Models: Implications for
Background and Pseudo-Absence Data.''} \emph{Ecological Applications} 19
(January): 181--97. \url{https://doi.org/10.1890/07-2153.1}.

\bibitem[\citeproctext]{ref-Pollack2012}
Pollack, Jennifer Beseres, Andrew Cleveland, Terence A. Palmer, Anthony
S. Reisinger, and Paul A. Montagna. 2012. {``A Restoration Suitability
Index Model for the Eastern Oyster (Crassostrea Virginica) in the
Mission-Aransas Estuary, TX, USA.''} \emph{PLoS ONE} 7 (July).
\url{https://doi.org/10.1371/journal.pone.0040839}.

\bibitem[\citeproctext]{ref-r}
R Core Team. 2024. \emph{R: A Language and Environment for Statistical
Computing}. Vienna, Austria: R Foundation for Statistical Computing.
\url{https://www.R-project.org/}.

\bibitem[\citeproctext]{ref-Roberts2017}
Roberts, David R., Volker Bahn, Simone Ciuti, Mark S. Boyce, Jane Elith,
Gurutzeta Guillera-Arroita, Severin Hauenstein, et al. 2017.
{``Cross-Validation Strategies for Data with Temporal, Spatial,
Hierarchical, or Phylogenetic Structure.''} \emph{Ecography}. Blackwell
Publishing Ltd. \url{https://doi.org/10.1111/ecog.02881}.

\bibitem[\citeproctext]{ref-Rushton2004}
Rushton, S. P., S. J. Ormerod, and G. Kerby. 2004. {``New Paradigms for
Modelling Species Distributions?''} \emph{Journal of Applied Ecology}.
\url{https://doi.org/10.1111/j.0021-8901.2004.00903.x}.

\bibitem[\citeproctext]{ref-Stachelek2015}
Stachelek, Joseph, and Christopher J. Madden. 2015. {``Application of
Inverse Path Distance Weighting for High-Density Spatial Mapping of
Coastal Water Quality Patterns.''} \emph{International Journal of
Geographical Information Science} 29 (7): 1240--50.
\url{https://doi.org/10.1080/13658816.2015.1018833}.

\bibitem[\citeproctext]{ref-TNC2021OysterRestoration}
The Nature Conservancy. 2021. {``Pensacola East Bay Oyster Habitat
Restoration Project.''} The Nature Conservancy.
\url{https://www.nature.org/content/dam/tnc/nature/en/documents/PEB-Oyster-Habitat-Restoration-Fact-Sheet-FINAL-7-13-21.pdf}.

\bibitem[\citeproctext]{ref-wqp_2026}
USGS, USDA, USEPA. 2026. {``Water Quality Portal.''} {U.S. Geological
Survey}; {U.S. Environmental Protection Agency}; {U.S. Department of
Agriculture}. \url{https://www.waterqualitydata.us/}.

\bibitem[\citeproctext]{ref-Valavi2019}
Valavi, Roozbeh, Jane Elith, José J. Lahoz-Monfort, and Gurutzeta
Guillera-Arroita. 2019. {``blockCV: An r Package for Generating
Spatially or Environmentally Separated Folds for k-Fold Cross-Validation
of Species Distribution Models.''} \emph{Methods in Ecology and
Evolution} 10 (February): 225--32.
\url{https://doi.org/10.1111/2041-210X.13107}.

\bibitem[\citeproctext]{ref-tidyverse}
Wickham, Hadley, Mara Averick, Jennifer Bryan, Winston Chang, Lucy
D'Agostino McGowan, Romain François, Garrett Grolemund, et al. 2019.
{``Welcome to the {tidyverse}.''} \emph{Journal of Open Source Software}
4 (43): 1686. \url{https://doi.org/10.21105/joss.01686}.

\end{CSLReferences}




\end{document}
